
\section{Semantic Compression}

\subsection{Entropy Compression vs Semantic Compression}
Let $X$ be a space of token sequences and $M$ a probabilistic language model. We distinguish two operators:
\begin{align*}
\text{entropy compression:} \quad &x \mapsto z_M(x), \\
\text{semantic compression:} \quad &x \mapsto k(x).
\end{align*}
Here $z_M(x)$ is a lossless binary code produced by entropy coding driven by model probabilities. Define model code length
\[
C_M(x) := |z_M(x)|,
\]
and a strong generic entropy-compressor baseline
\[
e(x),
\]
and per-token model-relative cross-entropy estimate
\[
h_M(x) := \frac{|z_M(x)|}{|x|_{\mathrm{tok}}}.
\]
LLM code length is a model-relative predictability signal, not a direct semantic observable. Coding efficiency ($C_M$, $e$, or $h_M$) is distinct from semantic invariance.

\subsection{Kolmogorov Complexity, Shannon Entropy, and Code-Length Rate}
These quantities are related but not identical.

For a random variable $X$ with distribution $p$, Shannon entropy is
\[
H(X) = -\sum_x p(x)\log_2 p(x).
\]
For autoregressive model probabilities $p_M(t_i\mid t_{<i})$, arithmetic coding yields code length approximately
\[
|z_M(x)| \approx -\sum_{i=1}^{n} \log_2 p_M(t_i\mid t_{<i}).
\]
Hence
\[
h_M(x) = \frac{|z_M(x)|}{|x|_{\mathrm{tok}}}
\approx
\frac{1}{n}\sum_{i=1}^{n} -\log_2 p_M(t_i\mid t_{<i}),
\]
which is an empirical model-relative cross-entropy rate (bits per token).

If $M$ matches the true source distribution, this rate approaches Shannon entropy rate; otherwise it is cross-entropy rate under model mismatch.

Kolmogorov complexity $K(x)$ is algorithmic (shortest program length) rather than probabilistic. Practical code lengths such as $|z_M(x)|$ and $e(x)$ are computable proxies for compressibility and may approximate $K(x)$ only asymptotically and model-relatively.

\subsection{NID and NCD}
For finite strings $x,y$, let $K(\cdot)$ denote prefix Kolmogorov complexity and $K(\cdot\mid\cdot)$ conditional complexity. Information distance is
\[
E(x,y) := \max\{K(x\mid y), K(y\mid x)\}.
\]
The normalized information distance (NID) is
\[
\mathrm{NID}(x,y) := \frac{E(x,y)}{\max\{K(x),K(y)\}}.
\]
NID is universal but incomputable.

Given a practical compressor $C$, normalized compression distance (NCD) is
\[
\mathrm{NCD}_C(x,y) :=
\frac{C(xy)-\min\{C(x),C(y)\}}{\max\{C(x),C(y)\}},
\]
where $xy$ denotes a self-delimiting concatenation.
Under standard normal-compressor assumptions, $\mathrm{NCD}_C$ approximates NID in practice.

In the NNCP-style setting, we instantiate
\[
C \equiv C_M, \quad C_M(x)=|z_M(x)|,
\]
which yields a model-relative distance
\[
\mathrm{NCD}_M(x,y) :=
\frac{|z_M(xy)|-\min\{|z_M(x)|,|z_M(y)|\}}{\max\{|z_M(x)|,|z_M(y)|\}}.
\]
Crucially, $z_M(x)$ is not a semantic projection: it is a codeword in bit space. A model can compress well because of semantics, but also because of boilerplate, repetition, formatting regularity, or domain familiarity.

To stabilize practical complexity estimates, we define a hybrid proxy
\[
C_M^*(x) := \min\{e(x),\,|z_M(x)|\}.
\]
This prevents a weak model predictor from degrading complexity estimates relative to a strong generic compressor and gives a more conservative practical proxy to Kolmogorov complexity.

We also define semantic surplus
\[
\Delta_M(x) := e(x)-|z_M(x)|.
\]
Large positive $\Delta_M(x)$ means $M$ predicts $x$ better than generic entropy compression; small or negative $\Delta_M(x)$ means little model advantage. Even large $\Delta_M$ is semantic evidence, not semantic proof.

Accordingly, compression distances are diagnostic rather than definitive. We use
\[
\mathrm{NCD}_{C_M^*}(x,y) :=
\frac{C_M^*(xy)-\min\{C_M^*(x),C_M^*(y)\}}{\max\{C_M^*(x),C_M^*(y)\}},
\]
and interpret them together with behavioral criteria.

\subsection{Limits of Entropy-Only Semantics}

Compression metrics such as $h_M(x)$, $\Delta_M(x)$, and $\mathrm{NCD}_{C_M^*}$ are useful diagnostics, but they do not define semantics. Low entropy does not guarantee semantic richness, and high entropy does not necessarily imply incoherence\footnote{High entropy may reflect semantic richness, noise, or unpredictability for other reasons; it is not a direct indicator of incoherence.}. Boilerplate, repetition, tautologies, domain familiarity, and even predictable nonsense can all yield low code lengths without semantic depth. Thus, compression alone cannot define semantics.

This motivates the need for behavioral semantic criteria: we must look beyond code length and compression distance to the actual behavior of the model on transformed inputs.

\subsection{Behavioral Semantic Preservation}
We now introduce the semantic compression operator
\[
k : X \to X
\]
which transforms text into a shorter or structurally simpler representation while attempting to preserve meaning. A motivating example is ``caveman compression,'' where grammatical structure and redundancy are aggressively removed:
\[
\texttt{"The quick brown fox jumps over the lazy dog"}
\;\mapsto\;
\texttt{"fox jump dog"}.
\]

Semantic preservation is defined behaviorally: if a model behaves similarly\footnote{Judging by the output similarity} on both $x$ and $k(x)$, then the compression preserves semantics relative to that model. Formally, fix a model $M$ and treat any evaluation context as implicit in $x$ itself. Let $f_M(x)$ denote model behavior on input $x$, and let $d_M$ be a behavioral distance on such outputs. We then define exact behavioral equivalence by
\[
x \sim_M y
\quad \Longleftrightarrow \quad
d_M\big(f_M(x), f_M(y)\big) = 0.
\]
Since exact equality is usually too rigid in practice, we use the relaxed relation
\[
x \approx_M^{(\varepsilon)} y
\quad \Longleftrightarrow \quad
d_M\big(f_M(x), f_M(y)\big) \le \varepsilon,
\]
for a chosen tolerance $\varepsilon > 0$. Semantic preservation under compression is then the statement
\[
x \approx_M^{(\varepsilon)} k(x)
\]
for the relevant evaluations.

Repeated application of $k$ may reveal stabilization:
\[
d_M\big(k(k(x)), k(x)\big) \le \varepsilon_{\mathrm{stab}}.
\]
This suggests a dynamical interpretation: repeated semantic compression can drive texts toward approximately invariant semantic regions.

To make this precise, let $d_M$ be a model-relative behavioral distance and let $\varepsilon > 0$. We define the approximate semantic kernel of $k$ (relative to $M$) as
\[
K_{k,M}^{(\varepsilon)} := \{x \in X : d_M(k(x),x) \le \varepsilon\}.
\]
Elements of $K_{k,M}^{(\varepsilon)}$ are approximately stable under semantic compression. A stronger set-level form is approximate invariance:
\[
k\big(K_{k,M}^{(\varepsilon)}\big) \subseteq K_{k,M}^{(\varepsilon')} 
\quad \text{for small } \varepsilon' \text{ (typically } \varepsilon' \approx \varepsilon\text{)}.
\]
Under this view, abstraction is interpreted as convergence of trajectories $x_{t+1}=k(x_t)$ toward these approximately invariant kernel regions, rather than exact fixed points.

\subsection{Open Hypotheses and Approach}
We now state open hypotheses, focusing on model-relative predictability and semantic stabilization, without invoking utilities prematurely.

\textbf{Hypothesis H1 (model-relative predictability regime):} For each model $M$, there exist calibrated thresholds
\[
\tau_{\mathrm{low}}(M) < \tau_{\mathrm{high}}(M)
\]
such that texts with $h_M(x) \in [\tau_{\mathrm{low}},\tau_{\mathrm{high}}]$ are more likely to preserve semantics (as measured by behavioral invariance) than texts outside this range. These thresholds describe predictability under $M$, not universal semantic quality.

\textbf{Tau-filtered preservation:} We say $x \sim_{\tau, M}^{(\varepsilon)} k(x)$ if $h_M(x)$ lies within the calibrated regime and
\[
x \approx_M^{(\varepsilon)} k(x)
\]
for the relevant evaluations.

Compression-based metrics such as $\mathrm{NCD}_{C_M^*}$ are auxiliary evidence, not primary semantic definitions.

\textbf{Operational approach:}
\begin{enumerate}
  \item Calibrate $\tau_{\mathrm{low}},\tau_{\mathrm{high}}$ from empirical percentiles of $h_M$ on mixed corpora.
  \item Measure $\Delta_M(x)=e(x)-|z_M(x)|$ to estimate model predictive advantage beyond generic compression.
  \item Measure $\mathrm{NCD}_{C_M^*}(x,k(x))$ and task-level output divergence as complementary semantic-preservation signals.
  \item Quantify stabilization by an idempotence residual such as
  \[
  r(x) := d_M\big(k(x), k(k(x))\big),
  \]
  where $d_M$ is a model-relative behavioral distance.
  \item Validate across diverse domains (encyclopedic text, scientific text, legal text, multilingual text, and code) to test cross-domain robustness.
\end{enumerate}

% --- Placeholders for further subsections ---
\subsection{Semantic Utilities and Utility-Conditioned Semantics (to be developed)}
% This section will introduce utility families as refinements of behavioral invariance, following the narrative order in intro.tex.

\subsection{Caveman Compression Example (to be developed)}
% This section will expand on the caveman example, closely following the style and clarity of intro.tex.
